\documentclass[11pt]{article}
\usepackage[margin=0.85in]{geometry}
\usepackage{amsmath,amssymb,amsthm,mathtools}
\usepackage{booktabs,longtable,tabularx,array,enumitem,microtype,needspace}
\usepackage{xcolor}
\usepackage{hyperref}
\pdfstringdefDisableCommands{\def\ensuremath#1{#1}\def\geq{≥}}
\hypersetup{
  colorlinks=true,
  linkcolor=blue!55!black,
  citecolor=blue!55!black,
  urlcolor=blue!55!black,
  pdfauthor={Brandon Li},
  pdftitle={No binary [[14,3,d\ensuremath{\geq}5]] quantum stabilizer code exists},
  pdfsubject={A signed-shadow proof and exhaustive additive-code classification},
  pdfcreator={LaTeX with Tectonic}
}
\setlength{\parindent}{0pt}
\setlength{\parskip}{4pt plus 1pt minus 1pt}
\setlength{\emergencystretch}{1.5em}
\setcounter{secnumdepth}{3}
\newtheorem{theorem}{Theorem}
\newtheorem{lemma}[theorem]{Lemma}
\newtheorem*{corollary*}{Corollary}
\providecommand{\tightlist}{\setlength{\itemsep}{0pt}\setlength{\parskip}{0pt}}
\providecommand{\real}[1]{#1}
\title{No binary \(\left[\!\left[14,3,d\geq 5\right]\!\right]\) quantum stabilizer code exists\\[0.7em]
\large\mdseries A signed-shadow proof and exhaustive additive-code classification}
\author{Brandon Li}
\date{}
\begin{document}
\maketitle
\begin{abstract}
We prove that no binary qubit stabilizer code with parameters
\(\left[\!\left[14,3,d\ge 5\right]\!\right]\) exists, settling the
remaining distance gap for stabilizer codes encoding three qubits into
fourteen. The proof first applies a signed affine-shadow inequality and
an integrality congruence to restrict low-weight stabilizer words.
Split-shadow Farkas certificates then exclude weight-two words
and every configuration with three weight-four stabilizer words, leaving
a unique weight-four stabilizer word.

That word yields a ten-coordinate binary-additive code of size 1,024,
minimum distance at least five, and symplectic hull dimension four. We
reconstruct all 37 relevant equivalence classes by exhaustive
lengthening: 25 have the wrong hull dimension, and the other twelve
violate the necessary coset-capacity bound. Together with the known
\(\left[\!\left[14,3,4\right]\!\right]\) construction, this proves
\(d_{\max}(14,3)=4\) for binary stabilizer codes.
\end{abstract}

\section{Introduction and main result}\label{introduction-and-main-result}

\begin{theorem}[Nonexistence]\label{thm:main}
There is no binary qubit stabilizer code with parameters \(\left[\!\left[14,3,d\ge 5\right]\!\right]\).
\end{theorem}

\begin{corollary*}
The maximum distance of a binary qubit stabilizer code with \(n=14\) and \(k=3\) is \(d_{\max}(14,3)=4\).
\end{corollary*}
\begin{proof}
Theorem~\ref{thm:main} gives the upper bound, and the known
\(\left[\!\left[14,3,4\right]\!\right]\) construction gives the matching lower bound
\cite{grassl-table}.
\end{proof}

Grassl's quantum-code tables \cite{grassl-table} record a \(\left[\!\left[14,3,4\right]\!\right]\) construction and an upper bound of five for the distance of a binary \(\left[\!\left[14,3\right]\!\right]\) stabilizer code. Ball, Centelles, and Huber \cite{ball-centelles-huber} identified \(\left[\!\left[14,3,5\right]\!\right]\) in Research Problem 1 as the smallest parameter set for which the existence of a qubit stabilizer code was then unknown. Theorem~\ref{thm:main} resolves that question in the negative. The nonexistence proof for \(\left[\!\left[13,5,4\right]\!\right]\) by Bierbrauer, Fears, Marcugini, and Pambianco \cite{bierbrauer-et-al} provides a related precedent. Our argument combines signed weight enumerators \cite{rains-weight,rains-shadow} with the classification of binary-additive codes; it does not assume \(\operatorname{GF}(4)\)-linearity.

The proof has three steps. First, signed-shadow identities restrict the possible low-weight stabilizer words. Second, linear infeasibility certificates eliminate the remaining low-weight configurations. Third, a symplectic reduction leads to an exhaustive classification of ten-site additive codes and a final coset-counting contradiction. The computational verification is summarized in Section~\ref{computational-material}.

Here, \emph{physical} refers to weight or support on the original qubits, as distinct from a coordinate of an auxiliary relaxation. The result applies to both pure and degenerate binary stabilizer codes, but not to arbitrary non-stabilizer quantum codes.

\section{Pauli space and notation}\label{pauli-space-and-notation}

Let

\[
\begin{aligned}
V_{n} &=\mathbb{F}_{2}^{2n} = \{(x\mid z): x,z \in \mathbb{F}_{2}^{n}\}
\end{aligned}
\]

encode Pauli labels modulo phase. The symplectic form is

\[
\begin{aligned}
\langle  (x\mid z),(x'\mid z')  \rangle &=x \cdot z' + z \cdot x'  \in \mathbb{F}_{2}.
\end{aligned}
\]

The Pauli weight is

\[
\begin{aligned}
\operatorname{wt}(x\mid z) &=\#\{i : (x_{i},z_{i}) \ne  (0,0)\}.
\end{aligned}
\]

A binary stabilizer code with parameters \(\left[\!\left[14,3,d\right]\!\right]\) is an isotropic subspace \(S  \le  V_{14}\) of binary dimension \(14-3=11\). Its centralizer is \(S^\perp\), and its distance is

\[
\begin{aligned}
d(S) &=\min \{ \operatorname{wt}(v) : v \in S^\perp \setminus  S \}.
\end{aligned}
\]

We call each element of \(S\) a stabilizer word. A stabilizer generator is an element of a chosen basis of \(S\); weight enumerators count all stabilizer words, not only basis generators.

For any nonzero binary subspace \(C\) of a Pauli space, write
\[
d(C)=\min\{\operatorname{wt}(c):c\in C\setminus\{0\}\}.
\]
We write \(d_{\max}(n,k)\) for the greatest minimum distance \(d(S)\) among binary stabilizer codes with fixed \(n\) and \(k\); this should not be confused with the maximum of the minimum weights of distinct logical-Pauli cosets within one fixed code.

We assume for contradiction that \(d(S) \ge 5\). A code is \emph{pure to distance d} if no nonzero stabilizer word has weight below \(d\); otherwise it is degenerate (or impure). In particular, the distance formula does not exclude low-weight stabilizer words. If a vector of weight at most four commutes with \(S\), it must itself be in \(S\).

For a binary subspace \(C\) of a Pauli space, its symplectic hull is \(\operatorname{Hull}(C)=C \cap C^\perp\). When a block of coordinates is specified, \(H_{a,b}\) below counts stabilizer words of split weight \((a,b)\); it is an enumerator, not a subspace. The hull in Section 6 is denoted \(J\) to keep the two notions distinct.

\section{The signed affine shadow, including split weights}\label{the-signed-affine-shadow-including-split-weights}

For one qubit define

\[
\begin{aligned}
q(x,z) &=x + z + xz \in \mathbb{F}_{2}.
\end{aligned}
\]

This is one exactly when \((x,z)\) is nonzero, so on \(V_{n}\) it satisfies

\[
\begin{aligned}
q(v) &=\operatorname{wt}(v) \pmod 2, \\
q(v+w) &=q(v) + q(w) + \langle v,w \rangle.
\end{aligned}
\]

Because \(S\) is isotropic, \(q\mid _{S}\) is linear. Nondegeneracy of the symplectic form gives a vector \(v_{0}\) such that

\[
\begin{aligned}
\langle v_{0},s \rangle &=q(s)\quad\forall s \in S.
\end{aligned}
\]

Character orthogonality gives the indicator identity

\[
\begin{aligned}
1_{v_{0}+S^\perp}(v) \\
& = 2^{-11} \sum_{s \in S} (-1)^{q(s)+\langle v,s \rangle}.
\end{aligned}
\]

Thus the signed transform counts weights in the affine coset \(v_{0}+S^\perp\), so every split cell is a nonnegative integer.

Fix a partition of the fourteen qubits into blocks of lengths \(m\) and \(f\), with \(m+f=14\), and let \(H_{i,j}\) count stabilizer words with block weights \((i,j)\). Define the quaternary Krawtchouk coefficients

\[
\begin{aligned}
K_{a}^{m}(i) &=[u^{a}] (1+3u)^{m-i} (1-u)^{i}.
\end{aligned}
\]

If \(C_{a,b}\) counts \(S^\perp\) and \(Sh_{a,b}\) counts the affine shadow, then

\[
\begin{aligned}
2048 C_{a,b} \\
& = \sum_{i,j} H_{i,j} K_{a}^{m}(i) K_{b}^{f}(j), \\
\\[0.4ex]
2048 Sh_{a,b} \\
& = \sum_{i,j} (-1)^{i+j} H_{i,j} \\
& K_{a}^{m}(i) K_{b}^{f}(j).
\end{aligned}
\]

For use in the split-weight systems, write \(T_{a,b}=2048C_{a,b}\) and \(\mathtt{SH}_{a,b}=2048Sh_{a,b}\); these are the integer transform numerators in the displayed formulas.

Both \(C_{a,b}\) and \(Sh_{a,b}\) are nonnegative integers, and \(C_{a,b} \ge H_{a,b}\) because \(S  \le  S^\perp\). If \(1 \le a+b \le 4\), then

\[
\begin{aligned}
C_{a,b} &=H_{a,b},
\end{aligned}
\]

because every centralizer word in that weight range must be a stabilizer word when \(d(S) \ge 5\). This equality is valid even when \(S\) is impure.

\section{Global low-weight restrictions and parity}\label{global-low-weight-restrictions-and-parity}

Write \(A_{j}\) for the number of stabilizer words of total physical weight \(j\), and put

\[
\begin{aligned}
T_{j} &=\sum_{w} A_{w} K_{j}^{14}(w), \\
U_{j} &=\sum_{w} (-1)^{w} A_{w} K_{j}^{14}(w).
\end{aligned}
\]

Write \(Sh_j=\sum_{a+b=j}Sh_{a,b}\) for the total-weight count in the affine shadow. The quantities \(T_{j}/2048\) and \(U_{j}/2048\) are the actual total-weight counts in \(S^\perp\) and in the affine shadow. For the fixed partition of Section~3, coefficient convolution gives

\[
\begin{aligned}
\sum_{a+b=j}K_a^m(i)K_b^f(\ell)
&=[t^j](1+3t)^{m-i}(1-t)^i(1+3t)^{f-\ell}(1-t)^\ell\\
&=K_j^{14}(i+\ell).
\end{aligned}
\]

Summing the split-shadow transform of Section~3 over \(a+b=j\) therefore gives
\[
U_j=\sum_{a+b=j}\mathtt{SH}_{a,b}=2048Sh_j.
\]

Hence \(U_{j} \ge 0\), even without purity. Also \(A_{0}=1\), \(\sum_{j} A_{j}=2048\), and \(T_{j}=2048 A_{j}\) for \(j=1,2,3,4\).

\begin{lemma}[Low-weight restriction]\label{lem:low-weight}
If \(d(S)\geq5\), then \(A_1=A_3=0\), \(A_2\leq2\), and \(A_4\leq3\).
\end{lemma}
\begin{proof}

Expanding the Krawtchouk transforms in \(T_{j}\) and \(U_{j}\) gives the coefficientwise identity

\[
\begin{aligned}
989184 A_{4} &=-296 U_{2} -88 U_{4} -9584640 A_{0} +6309 \sum_{j} A_{j} \\
& -11685888 A_{1} -1345536 A_{2} -4681728 A_{3} \\
& +1914 (T_{1}-2048 A_{1}) +1318 (T_{2}-2048 A_{2}) \\
& +279 (T_{3}-2048 A_{3}) +161 (T_{4}-2048 A_{4}) \\
& -1081344 A_{5} -254976 A_{6} -15360 A_{7} -44032 A_{14}.
\end{aligned}
\]

Now \(A_0=1\), \(\sum_j A_j=2048\), and \(T_j=2048A_j\) for \(1\leq j\leq4\). Since \(U_2,U_4,A_5,A_6,A_7,A_{14}\geq0\), the identity implies

\[
\begin{aligned}
11685888 A_{1} +1345536 A_{2} +4681728 A_{3} +989184 A_{4}  &\le  3336192.
\end{aligned}
\]

Since the \(A_j\) are nonnegative integers, these coefficients imply

\[
\begin{aligned}
A_{1} &=A_{3} = 0,   A_{2}  \le  2,   A_{4}  \le  3.
\end{aligned}
\]

\end{proof}

\noindent\emph{Remark.} Section~5.3 excludes the remaining cases \(A_{2}\in\{1,2\}\).

\subsection{Integral affine-shadow parity}\label{integral-affine-shadow-parity}

The inequality alone does not determine whether \(A_{4}\) is zero. An integrality congruence supplies the missing condition. We choose the coefficients below so that expanding the MacWilliams and shadow equations leaves, modulo 4096, only the terms in \(A_{3}\) and \(A_{4}\). Let

\[
\begin{aligned}
E_{0} &=\sum_{w} A_{w}=2048, \\
E_{j} &=T_{j} -2048 A_{j},                     j=1,\dots,4, \\
F_{j} &=U_{j} -2048 Sh_{j},                    j=0,\dots,3.
\end{aligned}
\]

Take the integer combination

\[
\begin{aligned}
L &=E_{0} -2E_{1} -6E_{2} -17E_{3} -95E_{4} \\
& -20F_{0} -24F_{1} +1108F_{2} +246F_{3}.
\end{aligned}
\]

The equations give \(L=2048\). Writing \(c_j\) for the coefficient of \(A_j\), direct expansion gives

\[
\begin{aligned}
(c_0,\dots,c_{14})={}&(-4550656,-7348224,-1052672,-2623488,321536,\\
&\qquad -634880,221184,-28672,24576,24576,\\
&\qquad -45056,0,40960,-16384,-77824).
\end{aligned}
\]

Modulo 4096 every coefficient vanishes except those of \(A_{3}\) and \(A_{4}\), each congruent to 2048. The four shadow-variable coefficients are \((40960, 49152, -2269184, -503808)\), all divisible by 4096. Since \(L=2048\), reduction modulo 4096 gives
\[
A_{3}+A_{4}\equiv1\pmod 2.
\]

Together with \(A_{3}=0\) and \(0 \le A_{4} \le 3\), this gives

\[
\begin{aligned}
A_{4} \in \{1,3\}.
\end{aligned}
\]

\section{Excluding low-weight stabilizer configurations}\label{excluding-low-weight-stabilizer-configurations}

For each specified subgroup, we derive a relaxation satisfied by every hypothetical code containing it. A Farkas certificate excludes the remaining low-weight configurations when the relaxation is infeasible.

\subsection{A necessary coset-counting relaxation}\label{a-necessary-coset-counting-relaxation}

Let \(R  \le  S\) be an isotropic subgroup supported on \(m\) of the original qubits, with the remaining \(f=14-m\) sites forming a suffix. Every element of \(S\) commutes with \(R\), so its prefix lies in \(R^\perp\). Also, \(S\) is a union of complete \(R\)-cosets.

For each complete coset \(Q\) of \(R\) in \(R^\perp\), let

\[
\begin{aligned}
p(Q)=(p_{0},\dots,p_{m}), \\
  p_{a} &=\#\{q \in Q : \operatorname{wt}_{\mathrm{pref}}(q)=a\},\qquad 0\leq a\leq m.
\end{aligned}
\]

Here \(\operatorname{wt}_{\mathrm{pref}}\) is Pauli weight on the original prefix sites.
Let \(y_{p,b}\) be the number of complete \(R\)-cosets selected by \(S\), aggregated over all suffix Pauli labels of weight \(b\), and grouped only by the prefix pattern \(p\). Then \(y_{p,b} \ge 0\) and

\[
\begin{aligned}
H_{a,b} &=\sum_{p} p_{a} y_{p,b}.
\end{aligned}
\]

This is a necessary model for a real code. It deliberately forgets the actual prefix coset label, the suffix Pauli label, and all global linear compatibility among different suffixes. It therefore enlarges the feasible set. A contradiction in this relaxation excludes a physical code; feasibility would prove nothing.

For every such model we use the following row families:

\begin{itemize}
\tightlist
\item
  \(H_{0,0}=1\) and \(\sum_{a,b}H_{a,b}=2048\);
\item
  physical coset equalities or inequalities forced by the selected subgroup;
\item
  the low-weight distance equations \(T_{a,b}-2048H_{a,b}=0\) for \(1 \le a+b \le 4\);
\item
  centralizer inclusion \(2048H_{a,b}-T_{a,b} \le 0\) in every split cell;
\item
  centralizer nonnegativity \(-T_{a,b} \le 0\) in every split cell;
\item
  signed affine-shadow nonnegativity \(-\mathtt{SH}_{a,b} \le 0\) in every split cell; and
\item
  \(y \ge 0\).
\end{itemize}

The transforms \(T_{a,b}\) and \(\mathtt{SH}_{a,b}\) use Pauli weights on the original prefix and suffix sites. The coset-capacity bound of Section~6.2 is not used in these exclusions.

\subsection{Exact Farkas sign convention}\label{exact-farkas-sign-convention}

After choosing the variables for a particular relaxation and scaling by 2048, write the necessary system as

\[
\begin{aligned}
M z &=rhs,     G z  \le  0,     z  \ge  0.
\end{aligned}
\]

A Farkas certificate consists of nonnegative \(\lambda\), free-sign \(\nu\), and a coefficient vector \(c\) satisfying

\[
\begin{aligned}
\lambda  &\ge  0, \\
c &=\lambda G + \nu M, \qquad c\geq 0, \\
& \nu \cdot rhs < 0.
\end{aligned}
\]
The inequality \(c\geq0\) is entrywise.

For a putative feasible \(z\),

\[
\begin{aligned}
0  \le  c \cdot z &=\lambda \cdot (G z) + \nu \cdot (M z) \\
 &\le  \nu \cdot rhs < 0,
\end{aligned}
\]

which is impossible. We check each certificate using exact integer arithmetic. The constraints are derived from the stabilizer problem before applying Farkas' lemma; the certificate shows that the resulting linear system is infeasible.

\subsection{Excluding weight-two stabilizers}\label{excluding-weight-two-stabilizers}

Suppose \(p=Z_{0} Z_{1}\) belongs to \(S\). A one-qubit Clifford acts transitively on the three nonidentity Pauli labels, so local Clifford conjugations on sites 0 and 1 send any weight-two stabilizer word to this representative while preserving weight and the code parameters. In particular, Hadamards send \(X_0X_1\) to \(Z_0Z_1\). The eight two-site Pauli labels commuting with \(p\) split into four complete cosets of \(\langle p \rangle\) with two-site pattern

\[
\begin{aligned}
(1,0,1)&:1, \qquad (0,2,0):1, \qquad (0,0,2):2.
\end{aligned}
\]
The integers after the colons are the coset multiplicities.

For every fixed suffix Pauli label, complete-coset closure therefore gives

\[
\begin{aligned}
H_{2,b}  &\ge  H_{0,b},   b=0,\dots,12.
\end{aligned}
\]

The 39 variables are the split counts \(H_{a,b}\) of stabilizer words with prefix weight \(a\in\{0,1,2\}\) on sites 0 and 1 and suffix weight \(b\in\{0,\ldots,12\}\) on the remaining sites.

\Needspace{10\baselineskip}
Every hypothetical code satisfies the following constraints:

\begin{itemize}
\item \(H_{0,0}=1\), \(\sum_{a,b}H_{a,b}=2048\), and \(H_{a,b}\geq0\);
\item the 11 equalities \(T_{a,b}=2048H_{a,b}\) for \(1\leq a+b\leq4\);
\item the 13 coset inequalities \(H_{2,b}\geq H_{0,b}\); and
\item in each of the 39 split cells, centralizer inclusion \(T_{a,b}\geq2048H_{a,b}\), centralizer nonnegativity \(T_{a,b}\geq0\), and affine-shadow nonnegativity \(\mathtt{SH}_{a,b}\geq0\).
\end{itemize}

Together with the two normalization equalities, these give 13 equality rows. The inequalities give 130 rows, in addition to the 39 nonnegativity conditions on the variables. This relaxation omits the further coset-pattern relations in Section~5.1.

The Farkas certificate has 13 nonnegative inequality multipliers and 10 nonzero equality multipliers. The residual vector \(c\) has 17 positive and 22 zero entries and satisfies the 39 column identities \(c=\lambda G+\nu M\). Its normalization is

\[
\begin{aligned}
\nu \cdot rhs &=-2{,}571{,}108{,}352.
\end{aligned}
\]

Thus \(A_{2}=0\) for every hypothetical code. Section~\ref{computational-material} gives the reproduction instructions and identifies both verification scripts.

\subsection{Excluding overlapping low-weight stabilizer words}\label{excluding-overlapping-low-weight-checks}

The support argument in Section 5.5 uses the following exclusions. Each subgroup acts on the original qubits, and its normalizer and cosets are enumerated before the linear relaxation is formed.

\subsubsection{A weight-three subgroup}\label{a-weight-three-subgroup}

For \(R=\langle \mathtt{ZZZ} \rangle\) on three original sites, the 16 normalizer cosets have prefix patterns

\[
\begin{aligned}
(0,0,1,1)&:12, \\
(0,1,1,0)&:3, \\
(1,0,0,1)&:1.
\end{aligned}
\]
Here each integer is the number of normalizer cosets with the indicated prefix-weight pattern.

The Farkas certificate satisfies the sign conditions of Section~5.2 on all 36 columns and has \(\nu \cdot rhs=-22{,}528\). It also excludes a weight-three stabilizer word without invoking Lemma~\ref{lem:low-weight}.

\subsubsection{The four-site Bell subgroup}\label{the-four-site-bell-subgroup}

For \(R=\langle \mathtt{XXXX},\mathtt{ZZZZ} \rangle\) on four original sites, the 16 complete cosets have patterns

\[
\begin{aligned}
(1,0,0,0,3)&:1, \\
(0,0,2,0,2)&:9, \\
(0,0,0,4,0)&:6.
\end{aligned}
\]
The integers after the colons are the coset multiplicities.

For every suffix weight \(b\),

\[
\begin{aligned}
H_{1,b}=0, \\
& H_{4,b}=H_{2,b}+3H_{0,b}.
\end{aligned}
\]

\Needspace{4\baselineskip}
The Farkas certificate satisfies the sign conditions of Section~5.2 on all 55 columns and has

\[
\begin{aligned}
\nu \cdot rhs &=-320{,}017{,}816{,}092{,}672.
\end{aligned}
\]

It excludes a pair of weight-four Bell stabilizer words with the same support.

\subsubsection{The five-site overlap subgroup}\label{the-five-site-overlap-subgroup}

If two commuting weight-four words overlap in three sites and differ on two of those sites, local Clifford conjugation gives

\[
\begin{aligned}
p &=X_{0} X_{1} Z_{2} Z_{3} I_{4}, \\
q &=Z_{0} Z_{1} Z_{2} I_{3} Z_{4}.
\end{aligned}
\]

All three nonzero words in \(\langle p,q \rangle\) have weight four on five sites. The normalizer has 64 four-element cosets and nine prefix-weight patterns. The Farkas certificate satisfies the sign conditions of Section~5.2 on all 90 columns and has

\[
\begin{aligned}
\nu \cdot rhs &=-52{,}501{,}014{,}528.
\end{aligned}
\]

\subsubsection{The two six-site overlap subgroups}\label{the-two-six-site-overlap-subgroups}

For the same-letter two-site overlap, use

\[
\begin{aligned}
R &=\langle  \mathtt{ZZZZII}, \mathtt{ZZIIZZ}  \rangle.
\end{aligned}
\]

The six-site normalizer has 256 four-element cosets and 13 prefix patterns. The corresponding 117-variable system has a Farkas certificate satisfying the sign conditions of Section~5.2 on all 117 columns, with \(\nu \cdot rhs=-8{,}820{,}736\).

For the crossed two-site overlap, use

\[
\begin{aligned}
R &=\langle  \mathtt{ZZZZII}, \mathtt{XXIIXX}  \rangle.
\end{aligned}
\]

The nonzero subgroup words have weights \(4,4,6\). Here the normalizer has 14 prefix patterns; the resulting 126-variable system has a Farkas certificate satisfying the sign conditions of Section~5.2 on all 126 columns, with \(\nu \cdot rhs=-22{,}212{,}608\). We check both six-site exclusions separately, enumerating the complete cosets in each case.

\subsection{Support geometry of weight-four stabilizer words}\label{support-geometry-of-weight-four-checks}

Let \(p,q\) be distinct commuting weight-four stabilizer words. Let \(t\) be the number of common support sites and \(h\) the number of common sites on which their nonidentity Pauli letters differ. Commutation makes \(h\) even, and

\[
\begin{aligned}
\operatorname{wt}(p+q) &=8 - 2t + h.
\end{aligned}
\]

The preceding subgroup exclusions remove every case with \(t \ge 2\):

{\small
\setlength{\tabcolsep}{8pt}
\renewcommand{\arraystretch}{1.08}
\begin{center}
\begin{tabular}{@{}rrl@{}}
\toprule
\(t\) & \(h\) & obstruction \\
\midrule
4 & 2 & weight-two stabilizer \\
4 & 4 & four-site Bell subgroup \\
3 & 0 & weight-two stabilizer \\
3 & 2 & five-site overlap subgroup \\
2 & 0 & six-site same-letter triangle subgroup \\
2 & 2 & six-site crossed-overlap subgroup \\
\bottomrule
\end{tabular}
\end{center}
}

The omitted case \(t=4,h=0\) would give the same Pauli word twice, not two distinct stabilizer words. Thus two distinct weight-four stabilizer words meet in at most one qubit. If they meet in one, commutation forces the same local Pauli letter there. Local Clifford transformations therefore map all weight-four stabilizer words to Z-only words simultaneously. This does not imply that the full stabilizer is CSS.

\begin{minipage}{\linewidth}
If \(A_{4}=3\), the three weight-four stabilizer words are linearly independent: the sum of any two has weight six or eight, not four. Their support sets are pairwise disjoint or meet in one site. Up to site permutation and reordering, there are exactly five support patterns:

{\small
\setlength{\tabcolsep}{4pt}
\begin{tabularx}{\linewidth}{@{}>{\raggedright\arraybackslash}p{2.7cm}Xr>{\centering\arraybackslash}p{2.1cm}r@{}}
\toprule
type & word supports & union & pair-product weights & triple-product weight \\
\midrule
0 intersections & \(\{0,1,2,3\};\{4,5,6,7\};\{8,9,10,11\}\) & 12 & \(8,8,8\) & 12 \\
1 intersection & \(\{0,1,2,3\};\{3,4,5,6\};\{7,8,9,10\}\) & 11 & \(6,8,8\) & 10 \\
2 intersections & \(\{0,1,2,3\};\{3,4,5,6\};\{6,7,8,9\}\) & 10 & \(6,8,6\) & 8 \\
3 distinct intersections & \(\{0,1,2,3\};\{3,4,5,6\};\{2,4,7,8\}\) & 9 & \(6,6,6\) & 6 \\
one common site & \(\{0,1,2,3\};\{3,4,5,6\};\{3,7,8,9\}\) & 10 & \(6,6,6\) & 10 \\
\bottomrule
\end{tabularx}
}
\end{minipage}

To see that the list is complete, record which of the three pairs of supports intersect. There can be zero, one, two, or three intersecting pairs. When exactly two pairs intersect, their intersection sites must be distinct: if both were the same, that site would lie in all three supports and the third pair would intersect too. When all three pairs intersect, either their three intersection sites differ or all three supports share one site. These are precisely the five cases above.

\subsection{The four overlapping rank-three geometries}\label{the-four-overlapping-rank-three-geometries}

For each of the four non-disjoint rows in the table, let \(R\) be the rank-three subgroup generated by the three displayed Z-only words. Its cosets in the normalizer are grouped by prefix-weight histograms. The relaxations are summarized below:

In each branch, the system includes the global equations

\[
\begin{aligned}
\sum_{a+b=j} H_{a,b} &=0,   j=1,2,3, \\
\sum_{a+b=4} H_{a,b} &=3.
\end{aligned}
\]

The first three follow from the low-weight restrictions and the weight-two exclusion; the last is the branch assumption \(A_{4}=3\). They are necessary conditions on the hypothetical code, not consequences of subgroup containment alone.

{\small
\setlength{\tabcolsep}{3pt}
\begin{tabularx}{\linewidth}{@{}X>{\centering\arraybackslash}p{1.0cm}>{\centering\arraybackslash}p{1.4cm}>{\centering\arraybackslash}p{1.6cm}>{\centering\arraybackslash}p{1.9cm}>{\centering\arraybackslash}p{2.1cm}@{}}
\toprule
geometry & sites & patterns & variables & positive multipliers & \(\nu\cdot rhs\) \\
\midrule
three distinct intersections & 9+5 & 40 & 240 & 12 & \(-86{,}470{,}656\) \\
one common site & 10+4 & 27 & 135 & 13 & \(-9{,}031{,}680\) \\
two intersections & 10+4 & 60 & 300 & 14 & \(-20{,}484{,}096\) \\
one intersection & 11+3 & 74 & 296 & 14 & \(-61{,}440\) \\
\bottomrule
\end{tabularx}
}

For each row, the Farkas certificate has nonnegative inequality multipliers and nonnegative residuals in every column, with the negative normalization shown. Farkas' lemma then excludes the system. These systems also impose the global equations

\[
\begin{aligned}
A_{1}=A_{2}=A_{3}=0,   A_{4}=3.
\end{aligned}
\]

\Needspace{8\baselineskip}
\subsection{\texorpdfstring{The disjoint \(4+4+4+2\) geometry}{The disjoint 4+4+4+2 geometry}}\label{the-disjoint-4442-geometry}

The only possible \(A_{4}=3\) geometry left by Section 5.6 would be

\[
\begin{aligned}
R &=\langle  \mathtt{ZZZZIIIIIIIIII}, \mathtt{IIIIZZZZIIIIII}, \mathtt{IIIIIIIIZZZZII}  \rangle.
\end{aligned}
\]

Each four-site block has 64 cosets of \(\langle \mathtt{ZZZZ} \rangle\) in its even-X-parity normalizer. The six weight patterns below have multiplicities totaling \(1+4+3+24+24+8=64\):

{\def\LTcaptype{none} % do not increment counter
\begin{longtable}[]{@{}lr@{}}
\toprule\noalign{}
pattern \((p_{0},p_{1},p_{2},p_{3},p_{4})\) & multiplicity \\
\midrule\noalign{}
\endhead
\bottomrule\noalign{}
\endlastfoot
\((1,0,0,0,1)\) & 1 \\
\((0,1,0,1,0)\) & 4 \\
\((0,0,2,0,0)\) & 3 \\
\((0,0,1,0,1)\) & 24 \\
\((0,0,0,2,0)\) & 24 \\
\((0,0,0,0,2)\) & 8 \\
\end{longtable}
}

For each of the three blocks and the two untouched sites, let \(y_{p,q,r,b}\) count complete \(R\)-cosets grouped by those literal patterns and suffix weight \(b\). There are \(6*6*6*3=648\) nonnegative variables. The grouping forgets the actual coset labels, suffix letters, and global linearity, so it is a relaxation.

The four-block split enumerator has \(5^3\cdot3=375\) cells, indexed by the three four-site block weights and the weight on the two untouched sites, and the system has 70 equalities. A Farkas certificate using the signed-shadow nonnegativity constraints gives

\[
\begin{aligned}
\lambda  &\ge  0, \\
c &=\lambda G + \nu M  \ge  0, \\
\nu \cdot rhs &=-125{,}829{,}120.
\end{aligned}
\]

The equality rows include the same global weight equations used in Section 5.6: \(A_{1}=A_{2}=A_{3}=0\) and the branch assumption \(A_{4}=3\). They do not follow from containment of the disjoint subgroup alone. All 648 residual coefficients are nonnegative (403 positive and 245 zero), so the Farkas certificate excludes this final geometry as well.

\begin{lemma}[Unique weight-four word]\label{lem:unique-check}
Every binary \(\left[\!\left[14,3,d\geq5\right]\!\right]\) stabilizer code has exactly one stabilizer word of weight four.
\end{lemma}
\begin{proof}
The parity restriction gives \(A_4\in\{1,3\}\), and Sections~5.3--5.7 exclude \(A_4=3\). Hence \(A_4=1\).
\end{proof}

\section{Reduction to a ten-site additive code}\label{reduction-to-a-ten-site-additive-code}

\Needspace{11\baselineskip}
\begin{lemma}[Ten-site reduction]\label{lem:ten-site}
Suppose \(S\leq V_{14}\) is isotropic of dimension 11, has \(d(S)\geq5\), and its only nonzero stabilizer word of physical weight at most four is a word \(p\) of weight four. Equivalently, \(A_1=A_2=A_3=0\) and \(A_4=1\). Then there exists a binary subspace \(C\leq V_{10}\) of dimension 10, minimum Pauli weight at least five, and \(\dim(C\cap C^\perp)=4\).
\end{lemma}
\begin{proof}

Let \(p\) be the unique weight-four stabilizer word. After local Clifford conjugation and a site permutation, assume \(p=Z_{0}Z_{1}Z_{2}Z_{3}\) on the first four original qubits. Write \(V_{4}\) for their Pauli space and \(W_{10}\) for the remaining ten-qubit Pauli space, so \(V_{14}=V_{4} \oplus W_{10}\).

Define the physical six-dimensional quotient

\[
\begin{aligned}
E &=(p^\perp \cap V_{4}) / \langle p \rangle.
\end{aligned}
\]

The quotient is nondegenerate symplectic: \(p^\perp \cap V_{4}\) has binary dimension seven, and quotienting by its one-dimensional radical \(\langle p \rangle\) leaves dimension six.

By the hypothesis, \(S \cap V_{4}=\langle p \rangle\), so projection of \(S/\langle p \rangle\) to \(W_{10}\) is injective. Let its image be a binary ten-space \(P \le W_{10}\). Then there is a linear map \(T:P\to E\) such that

\[
\begin{aligned}
S/\langle p \rangle &=\{ T(x)+x : x \in P \}, \\
\langle x,y \rangle_{W} &=\langle T(x),T(y) \rangle_{E}.
\end{aligned}
\]

\subsection{Distance forces surjectivity}\label{distance-forces-surjectivity}

Suppose \(z \in E\) is orthogonal to \(\operatorname{im}(T)\). Choose a four-site representative \(v\) of \(z\) in \(p^\perp\). It commutes with \(p\) and with every element of the graph, hence \(v \in S^\perp\). Its weight on the original qubits is at most four. Distance at least five forces \(v \in S\). But \(S \cap V_{4}=\langle p \rangle\), so \(z=0\). Thus \((\operatorname{im}(T))^\perp=0\). Nondegeneracy of \(E\) gives \(\operatorname{im}(T)=E\), and therefore

\[
\begin{aligned}
\operatorname{rank}(T)=6,   \dim \operatorname{ker}(T)=10-6=4.
\end{aligned}
\]

Surjectivity also identifies the radical of the pairing on \(P\):

\[
\begin{aligned}
\operatorname{rad}(P)=P \cap P^\perp_{W}=\operatorname{ker}(T),   \dim \operatorname{rad}(P)=4.
\end{aligned}
\]

Put

\[
\begin{aligned}
C &=P^\perp_{W}.
\end{aligned}
\]

Then \(\dim C=10\), so \(\lvert C\rvert=1024\). If a nonzero \(c \in C\) had physical weight at most four, the suffix-only vector \((0,c)\) would be in \(S^\perp\), and distance would force it into \(S\). It cannot equal \(p\), since its prefix is zero; it would therefore be another nonzero stabilizer word of weight at most four, contrary to the hypothesis. Thus \(d(C) \ge 5\). Because \(C^\perp=P\),

\[
\begin{aligned}
\dim(C \cap C^\perp_{W})=\dim(P \cap C)=4.
\end{aligned}
\]

\end{proof}

\subsection{A necessary coset-capacity bound}\label{a-necessary-coset-capacity-bound}

Let \(J=C \cap C^\perp_{W}\), so \(\dim J=4\), and let \(P=C^\perp_{W}\). The quotient \(P/J\) is a nondegenerate six-dimensional symplectic space. Also \(J^\perp_{W}/C\) has 64 suffix cosets. For \(\xi=w+C\in J^\perp_{W}/C\), put
\[
d(\xi)=\min_{c\in C}\operatorname{wt}(w+c).
\]
\begin{minipage}{\linewidth}
The pairing
\[
(J^\perp_{W}/C)\times(P/J)\longrightarrow\mathbb{F}_{2},
\qquad
(w+C,x+J)\longmapsto\langle w,x\rangle_{W}
\]
is well defined and nondegenerate: its left kernel is \(P^\perp_{W}=C\), and both quotient spaces have dimension six.
\end{minipage}

For each \(\xi=w+C\), there is a unique \(e(\xi)\in E\) such that
\[
\langle e(\xi),T(x)\rangle_{E}=\langle w,x\rangle_{W}
\qquad(x\in P).
\]
The symplectic form on \(P/J\) identifies \((P/J)^*\) with \(P/J\), and \(T\) induces a symplectic isomorphism \(P/J\to E\). The perfect pairing above therefore identifies \(J^\perp_{W}/C\) with \(E\), so \(\xi\mapsto e(\xi)\) is a bijection and sends zero to zero.

We use a Hall-type counting argument to bound how many low-cost prefix classes a valid code can use. For \(e\in E\), define its physical prefix cost by
\[
\operatorname{cost}(e)=
\min\{\operatorname{wt}(v):v\in p^\perp\cap V_{4},\ v+\langle p\rangle=e\}.
\]
Its distribution is

\[
\begin{aligned}
\begin{array}{c|ccccc}
\text{cost} & 0 & 1 & 2 & 3 & 4\\
\hline
\text{number of classes} & 1 & 4 & 27 & 24 & 8
\end{array}
\end{aligned}
\]

To obtain these counts, split the even-weight \(x\)-masks in \(p^\perp\) by weight. The \(x=0\) quotient contributes \(1,4,3\) classes of costs \(0,1,2\); each of the six weight-two \(x\)-masks contributes four classes of cost \(2\) and four of cost \(3\); the weight-four mask contributes eight classes of cost \(4\).

For every suffix representative \(w\in\xi\), a prefix representative \(v\) of \(e(\xi)\) makes \(v+w\) commute with \(S\). Conversely, commutation forces this prefix class for a given suffix class. Thus the mixed centralizer fiber indexed by \(\xi\) has minimum weight \(d(\xi)+\operatorname{cost}(e(\xi))\). If \(\xi\ne0\), its suffix is nonzero, so no vector in this fiber is \(p\). Each vector in the fiber therefore has weight at least five: if it lies in \(S\), Sections~4--5 leave \(p\) as the only nonzero stabilizer word of weight at most four; otherwise its weight is at least \(d(S)\). Hence
\[
d(\xi)+\operatorname{cost}(e(\xi))\ge5
\qquad(\xi\ne0).
\]

Every compatible graph map induces a symplectic isomorphism \(P/J\to E\). The bijection \(\xi\mapsto e(\xi)\) and the cost distribution give the following capacities without enumerating these isomorphisms:

\[
\begin{aligned}
\#\{\xi \ne  0 : d(\xi) &\le 1\}  \le  8, \\
\#\{\xi \ne  0 : d(\xi) &\le 2\}  \le  32, \\
\#\{\xi \ne  0 : d(\xi) &\le 3\}  \le  59.
\end{aligned}
\]

The capacities are the numbers of nonzero prefix classes of cost at least 4, 3, and 2, respectively. Thus every symplectic lift must satisfy the three inequalities; failure of any one excludes it.

\section{Exhaustive additive-code classification}\label{exhaustive-additive-code-classification}

\subsection{What is classified}\label{what-is-classified}

Under the standard Pauli/additive-code correspondence \cite{calderbank-shor-sloane}, represent a quaternary symbol by two binary bits. A length-\(n\) additive code is a binary subspace of \(\mathbb{F}_{2}^{2n}\), with quaternary Hamming weight equal to physical Pauli weight. Equivalence is coordinate permutation plus an arbitrary local \(\operatorname{GL}(2,2)\) map on each two-bit symbol plane. This is exactly the binary local symplectic monomial equivalence relevant here. No global \(\operatorname{GF}(4)\) scalar-linearity is assumed.

Grassl, Krotov, Sok, and Solé \cite{grassl-et-al} report 37 equivalence classes in the \((n,k)=(10,10)\) cell of their additive \((n,2^{k}, \ge 5)_{4}\) classification (Section 2.3, Table 1 of \emph{The punctured dodecacode is unique}). They report the count, not the 37 generator matrices. We reconstruct the classes by lengthening. The published counts serve as cross-checks: the census generates candidates from predecessor cells and decides equivalence by the graph certificates described below; the table counts do not by themselves establish completeness.

\subsection{Why the lengthening census is exhaustive for this cell}\label{why-the-lengthening-census-is-exhaustive-for-this-cell}

The census builds additive codes one coordinate at a time. The projection argument shows that every code in the target class arises from one of the extension cases below.

Let \(D\leq\mathbb{F}_2^{2(m+1)}\) be a target code of binary dimension \(k\), and project it onto its last two-bit coordinate. The image has binary dimension \(r\in\{0,1,2\}\); the kernel is a length-\(m\) parent code of binary dimension \(k-r\). Choose a basis of the image and lift its basis vectors to \(D\). Their old-coordinate parts determine \(r\) quotient classes modulo the parent. For \(r=1\), the nonzero class must have minimum old weight at least four. For \(r=2\), each of the three nonzero combinations of the two classes must have minimum old weight at least four. After a local \(\operatorname{GL}(2,2)\) change of basis, these are exactly the extensions below. Conversely, each listed extension has distance at least five whenever its parent has distance at least five and the stated coset minima are at least four.

Thus lengthening a parent of length \(m\) and binary dimension \(k-r\) uses one new coordinate of rank \(r\):

\begin{itemize}
\tightlist
\item
  \(r=0\): append a zero coordinate;
\item
  \(r=1\): choose one quotient coset of the parent whose old minimum weight is at least four and attach a nonzero new symbol;
\item
  \(r=2\): choose two quotient labels whose three nonzero XOR combinations all have old minimum weight at least four and attach independent new symbols.
\end{itemize}

The threshold is four: a word with nonzero new symbol gains one physical unit of weight, while a word with zero new symbol lies in the distance-five parent. Thus every target code appears in one of these three cases.

\Needspace{12\baselineskip}
The table lists the predecessor cells needed to reach \((10,10)\) and the reconstructed class counts:

{\def\LTcaptype{none} % do not increment counter
\begin{longtable}[]{@{}rl@{}}
\toprule\noalign{}
length \(n\) & required binary dimensions \(k\): number of classes \\
\midrule\noalign{}
\endhead
\bottomrule\noalign{}
\endlastfoot
5 & \(0:1, 1:1, 2:1, 3:0, 4:0\) \\
6 & \(2:5, 3:1, 4:0, 5:0, 6:0\) \\
7 & \(4:43, 5:1, 6:0, 7:0, 8:0\) \\
8 & \(6:1579, 7:0, 8:0, 9:0, 10:0\) \\
9 & \(8:2298, 9:0, 10:0\) \\
10 & \(10:37\) \\
\end{longtable}
}

Every zero-count predecessor cell listed above is computed and verified empty rather than treated as an absent cell. Hence the table covers every predecessor needed by the recurrence through \((10,10)\).

\subsection{Equivalence certificates}\label{equivalence-certificates}

For every candidate the census constructs a colored incidence graph with one vertex for every codeword, three nonzero-symbol vertices for every coordinate, and one coordinate-group vertex binding each triple of symbols. The colors separate the three vertex types. A color-preserving graph isomorphism is exactly a reordering of codewords, a coordinate permutation, and an arbitrary permutation of the three nonzero symbols in each coordinate. The latter is \(\operatorname{GL}(2,2)\). Thus the graph certificate gives a canonical additive-monomial equivalence test; the weight and pair-rank invariants only filter candidates before canonicalization. The graph-isomorphism implementation is based on nauty \cite{mckay-piperno}.

\subsection{Hull and coset-capacity results}\label{hull-and-coset-capacity-results}

The census produces 37 generator matrices, each defining 1,024 words of minimum distance five. Separate audits recompute their ranks, distances, hull dimensions, and coset minima. The hull dimensions are

{\def\LTcaptype{none} % do not increment counter
\begin{longtable}[]{@{}rr@{}}
\toprule\noalign{}
binary symplectic hull dimension & number of classes \\
\midrule\noalign{}
\endhead
\bottomrule\noalign{}
\endlastfoot
0 & 7 \\
2 & 14 \\
4 & 12 \\
6 & 2 \\
8 & 2 \\
\end{longtable}
}

\Needspace{17\baselineskip}
The 25 classes with hull dimension different from four cannot arise from the code \(C\) required by Lemma~\ref{lem:ten-site}. For the 12 hull-four classes, the coset-capacity audit gives:

{\small
\setlength{\tabcolsep}{4pt}
\begin{tabularx}{\linewidth}{@{}rXrrrl@{}}
\toprule
class & suffix weight spectrum \(\{w:N_w\}\) & \(N_{\leq1}\) & \(N_{\leq2}\) & \(N_{\leq3}\) & capacity \\
\midrule
9 & \(\{0:1,2:39,3:24\}\) & 0 & 39 & 63 & fails \\
12 & \(\{0:1,1:2,2:37,3:24\}\) & 2 & 39 & 63 & fails \\
14 & \(\{0:1,2:39,3:24\}\) & 0 & 39 & 63 & fails \\
15 & \(\{0:1,2:27,3:36\}\) & 0 & 27 & 63 & fails \\
17 & \(\{0:1,1:2,2:37,3:24\}\) & 2 & 39 & 63 & fails \\
20 & \(\{0:1,2:35,3:27,4:1\}\) & 0 & 35 & 62 & fails \\
21 & \(\{0:1,1:2,2:37,3:24\}\) & 2 & 39 & 63 & fails \\
22 & \(\{0:1,2:23,3:40\}\) & 0 & 23 & 63 & fails \\
26 & \(\{0:1,2:35,3:27,4:1\}\) & 0 & 35 & 62 & fails \\
34 & \(\{0:1,1:2,2:37,3:24\}\) & 2 & 39 & 63 & fails \\
35 & \(\{0:1,1:2,2:37,3:24\}\) & 2 & 39 & 63 & fails \\
37 & \(\{0:1,1:2,2:37,3:24\}\) & 2 & 39 & 63 & fails \\
\bottomrule
\end{tabularx}
}

In every row \(N_{\leq3}\in\{62,63\}>59\). Eight rows also have \(N_{\leq2}=39>32\).

\begin{lemma}[Additive lift exclusion]\label{lem:additive-lift}
None of the 37 length-ten classes can supply the code in Lemma~\ref{lem:ten-site}: 25 have a different hull dimension, and the other twelve violate the coset-capacity bound.
\end{lemma}
\begin{proof}
The first exclusion follows from the hull-dimension table. For every remaining class, the coset table gives \(N_{\leq3}\in\{62,63\}\), exceeding the necessary capacity 59. Hence none of the 37 classes can supply the required code.
\end{proof}

\section{Discussion and conclusion}\label{discussion-and-conclusion}

\begin{proof}[Proof of Theorem~\ref{thm:main}]
Assume a binary \(\left[\!\left[14,3,d\ge 5\right]\!\right]\) stabilizer code exists. Sections~3--4 give \(A_{1}=A_{3}=0\), \(A_{2}\leq2\), and \(A_{4}\in\{1,3\}\). Section~5.3 then excludes weight-two stabilizer words, so \(A_{2}=0\). Sections~5.3--5.5 exclude all pairs of distinct weight-four stabilizer words whose supports intersect in at least two sites and classify the \(A_4=3\) case into five support geometries. Section~5.6 excludes the four overlapping rank-three geometries, and Section~5.7 excludes the disjoint \(4+4+4+2\) geometry. Hence \(A_{4}=1\).

Lemma~\ref{lem:ten-site} forces a ten-site additive code of size 1,024, distance at least five, and hull dimension four. Section~7 exhausts the 37 equivalence classes of such ten-site additive codes and rejects every one: 25 by hull dimension and 12 by the coset-capacity bound. This is a contradiction.
\end{proof}

The proof applies to these parameters and to stabilizer codes; it gives no bound at other lengths or for non-stabilizer codes.

\section{Computational material}\label{computational-material}

The verification uses Farkas certificates for the split-weight systems and an exhaustive classification of length-ten additive codes. The standard replay requires Python 3.10 or newer and Node.js. From the repository root, run
\begin{center}
\texttt{bash experiments/qec1435\_replay\_nonexistence.sh}
\end{center}
This replay checks the certificate arithmetic and the committed class representatives but does not regenerate the classification. The optional \texttt{--census} mode reconstructs the classification and requires the pinned \texttt{pynauty==2.8.8.1} dependency. The pin and detailed verification instructions are included in the repository. Replay outputs are temporary and do not replace the committed representatives or audit data.

\Needspace{4\baselineskip}
The public repository is
\begin{center}
\url{https://github.com/brandonlign/quantum-code-bounds}
\end{center}


\Needspace{20\baselineskip}
\begin{thebibliography}{99}
\bibitem{calderbank-shor-sloane}
A.~R. Calderbank, E.~M. Rains, P.~W. Shor, and N.~J.~A. Sloane,
``Quantum error correction via codes over GF(4),''
\emph{IEEE Trans. Inf. Theory} 44 (1998), 1369--1387.
\href{https://doi.org/10.1109/18.681315}{doi:10.1109/18.681315}.

\bibitem{rains-weight}
E.~M. Rains, ``Quantum weight enumerators,''
\emph{IEEE Trans. Inf. Theory} 44 (1998), 1388--1394.
\href{https://doi.org/10.1109/18.681316}{doi:10.1109/18.681316}.

\bibitem{rains-shadow}
E.~M. Rains, ``Quantum shadow enumerators,''
\emph{IEEE Trans. Inf. Theory} 45 (1999), 2361--2366.
\href{https://doi.org/10.1109/18.796376}{doi:10.1109/18.796376}.

\bibitem{ball-centelles-huber}
S. Ball, A. Centelles, and F. Huber, ``Quantum error-correcting codes and their geometries,''
\emph{Ann. Inst. H. Poincaré D} 10 (2023), no.~2, 337--405.
\href{https://doi.org/10.4171/AIHPD/160}{doi:10.4171/AIHPD/160}.

\bibitem{bierbrauer-et-al}
J. Bierbrauer, R. Fears, S. Marcugini, and F. Pambianco,
``The nonexistence of a [[13,5,4]]-quantum stabilizer code,''
\emph{IEEE Trans. Inf. Theory} 57 (2011), 4788--4793.
\href{https://doi.org/10.1109/TIT.2011.2146430}{doi:10.1109/TIT.2011.2146430}.

\bibitem{grassl-et-al}
M. Grassl, D. Krotov, L. Sok, and P. Solé,
``The punctured dodecacode is unique,''
\emph{Mathematics and Education in Mathematics} 55 (2026), 393--404.
\href{https://doi.org/10.55630/mem.2026.55.393-404}{doi:10.55630/mem.2026.55.393-404}.

\bibitem{mckay-piperno}
B.~D. McKay and A. Piperno, ``Practical graph isomorphism, II,''
\emph{J. Symbolic Comput.} 60 (2014), 94--112.
\href{https://doi.org/10.1016/j.jsc.2013.09.003}{doi:10.1016/j.jsc.2013.09.003}.

\bibitem{grassl-table}
M. Grassl, ``Bounds on the minimum distance of additive quantum codes, (n=14,k=3,q=4),''
maintained code table, accessed 24 September 2026.
\href{https://codetables.de/QECC.php?k=3&n=14&q=4}{codetables.de}.
\end{thebibliography}

\end{document}
